% ============================================================
%  Chapter 4 — System Design
% ============================================================
\chapter{System Design}
\label{cha:system_design}

% TODO: write full chapter — skeleton below
% Target: 14–16 pages

% ------------------------------------------------------------
\section{Design Principles}
\label{sec:design_principles}

% TODO: 1 page
% - Strategy pattern: chunking / embedding / RAG mode configurable per project at runtime
% - Event-driven ingestion (async) vs. synchronous query pipeline
% - Non-root Docker, structured JSON logging, distributed trace IDs (X-Trace-Id)
% - No module-level side effects; dependency injection via FastAPI Depends

% ------------------------------------------------------------
\section{System Context (C4 Level 1)}
\label{sec:c4_context}

% TODO: 1 page + figure

\begin{figure}[htbp]
\centering
\begin{tikzpicture}[
  actor/.style  = {rectangle, rounded corners=3pt,
                   draw=gray!60, fill=gray!8, thick,
                   minimum width=2.8cm, minimum height=0.85cm, align=center,
                   font=\small},
  system/.style = {rectangle, rounded corners=8pt,
                   draw=blue!55, fill=blue!7, very thick,
                   minimum width=5.5cm, minimum height=3.2cm, align=center},
  ext/.style    = {rectangle, rounded corners=4pt,
                   draw=purple!50, fill=purple!5, thick,
                   dashed, minimum width=3.0cm, minimum height=0.85cm,
                   align=center, font=\small},
  arr/.style    = {-{Stealth[length=5pt]}, thick, black!60},
  lbl/.style    = {font=\footnotesize, fill=white, inner sep=1pt},
]

% actors
\node[actor] (user)   at (-5.5, 0) {Patient\,/\,User\\[-1pt]{\tiny queries drug interactions}};
\node[actor] (admin)  at (-5.5,-2.0) {Administrator\\[-1pt]{\tiny manages projects \& docs}};

% system box
\node[system] (medrag) at (0, -1.0) {};
\node[font=\small\bfseries, blue!60] at (0, 0.6)  {\medrag{} Platform};
\node[font=\footnotesize]            at (0, -0.1) {14 microservices};
\node[font=\footnotesize]            at (0, -0.7) {FastAPI + Weaviate + MongoDB};
\node[font=\footnotesize]            at (0, -1.3) {RabbitMQ ingestion pipeline};

% external dependencies
\node[ext] (llm)     at (5.5,  0.2) {Anthropic LLM\\[-1pt]{\tiny Claude API}};
\node[ext] (embedpv) at (5.5, -2.2) {Embedding Provider\\[-1pt]{\tiny BGE-M3 / Cohere / OpenAI}};

% arrows
\draw[arr] (user.east)   -- node[lbl,above=1pt]{REST / SSE}       (medrag.west |- user);
\draw[arr] (admin.east)  -- node[lbl,above=1pt]{REST}             (medrag.west |- admin);
\draw[arr] (medrag.east |- llm)     -- node[lbl,above=1pt]{HTTPS} (llm.west);
\draw[arr] (medrag.east |- embedpv) -- node[lbl,above=1pt]{HTTPS} (embedpv.west);

\end{tikzpicture}
\caption{System context diagram. Two actors (Patient/User and Administrator)
  interact with the \medrag{} system boundary via REST and SSE. External
  dependencies are the LLM provider (Anthropic Claude) and the configurable
  embedding provider (BGE-M3 by default).}
\label{fig:c4_context}
\end{figure}

% ------------------------------------------------------------
\section{Container Architecture (C4 Level 2)}
\label{sec:c4_container}

% TODO: 3 pages + large figure

The \medrag{} platform consists of 14 services, listed in Table~\ref{tab:services}.

\begin{table}[htbp]
  \centering
  \caption{Service inventory: all containers, ports, and primary responsibilities.}
  \label{tab:services}
  \begin{tabularx}{\textwidth}{@{}llX@{}}
    \toprule
    \textbf{Service} & \textbf{Port} & \textbf{Responsibility} \\
    \midrule
    api-gateway     & 8000 & JWT validation, request routing to downstream services \\
    auth            & 8001 & User registration, login, JWT issuance (PyJWT + bcrypt) \\
    orchestrator    & 8002 & Query flow coordination, RAG mode dispatch, conversation storage \\
    query-processor & 8003 & Query rewriting and HyDE document generation \\
    retrieval       & 8004 & Hybrid search (Weaviate BM25 + vector, configurable alpha) \\
    reranker        & 8005 & Cross-encoder reranking (BGE-reranker-v2-m3) \\
    generation      & 8006 & LLM call, SSE streaming, citation extraction \\
    ingestion       & 8007 & File upload, content-hash deduplication, MQ publish \\
    parser          & 8008 & PDF and document text extraction (pypdf / docling) \\
    chunking        & 8009 & Configurable chunking: fixed-512, recursive, semantic \\
    embedding       & 8010 & Text embedding: local BGE-M3, Cohere, or OpenAI \\
    indexing        & 8011 & Weaviate vector insert and document status update \\
    admin           & 8012 & Project CRUD, document listing, re-index trigger \\
    eval            & 8013 & RAGAS metrics, async consumer of \texttt{query.completed} events \\
    \bottomrule
  \end{tabularx}
\end{table}

\begin{figure}[p]
\centering
\begin{tikzpicture}[
  svc/.style   = {rectangle, rounded corners=3pt,
                  draw=blue!50, fill=blue!6, thick,
                  minimum width=3.0cm, minimum height=0.65cm,
                  align=center, font=\footnotesize},
  mq/.style    = {rectangle, rounded corners=3pt,
                  draw=orange!60, fill=orange!8, thick,
                  minimum width=3.0cm, minimum height=0.65cm,
                  align=center, font=\footnotesize},
  infra/.style = {rectangle, rounded corners=3pt,
                  draw=gray!60, fill=gray!10, thick,
                  minimum width=2.6cm, minimum height=0.65cm,
                  align=center, font=\footnotesize},
  ext/.style   = {rectangle, rounded corners=3pt,
                  draw=purple!50, fill=purple!5, thick, dashed,
                  minimum width=2.8cm, minimum height=0.65cm,
                  align=center, font=\footnotesize},
  actor/.style = {rectangle, rounded corners=3pt,
                  draw=gray!55, fill=gray!7, thick,
                  minimum width=2.4cm, minimum height=0.65cm,
                  align=center, font=\footnotesize},
  supp/.style  = {rectangle, rounded corners=3pt,
                  draw=teal!50, fill=teal!5, thick,
                  minimum width=3.0cm, minimum height=0.65cm,
                  align=center, font=\footnotesize},
  arr/.style   = {-{Stealth[length=4pt]}, thick, black!60},
  lbl/.style   = {font=\scriptsize, fill=white, inner sep=1pt},
]

% ── Clients (y = 9) ──────────────────────────────────────────────
\node[actor] (user)    at (2.5, 9.0) {User\\[-1pt]{\tiny REST / SSE}};
\node[actor] (admin_u) at (5.5, 9.0) {Admin User\\[-1pt]{\tiny REST}};

% ── Gateway tier (y = 7.5) ───────────────────────────────────────
\node[svc]   (gw)   at (2.5, 7.5) {api-gateway\\\tiny:8000};
\node[svc]   (auth) at (9.5, 7.5) {auth\\\tiny:8001};

% ── Query Pipeline group (x = 0–6, y = 2.2–6.5) ─────────────────
\draw[rounded corners=6pt, fill=blue!3, draw=blue!35, dashed]
    (-0.2, 2.0) rectangle (5.7, 6.8);
\node[font=\footnotesize\bfseries, blue!50] at (2.75, 6.5) {Query Pipeline};

\node[svc] (orch)   at (2.5, 5.9) {orchestrator\\\tiny:8002};
\node[svc] (qp)     at (2.5, 5.1) {query-processor\\\tiny:8003};
\node[svc] (retr)   at (2.5, 4.3) {retrieval\\\tiny:8004};
\node[svc] (rerank) at (2.5, 3.5) {reranker\\\tiny:8005};
\node[svc] (gen_s)  at (2.5, 2.7) {generation\\\tiny:8006};

% ── Ingestion Pipeline group (x = 6.5–13, y = 1.4–6.5) ──────────
\draw[rounded corners=6pt, fill=orange!3, draw=orange!40, dashed]
    (6.5, 1.2) rectangle (12.7, 6.8);
\node[font=\footnotesize\bfseries, orange!60] at (9.6, 6.5) {Ingestion Pipeline};

\node[svc]   (ingest) at (9.5, 5.9) {ingestion\\\tiny:8007};
\node[mq]    (rmq_s)  at (9.5, 5.1) {RabbitMQ broker\\\tiny async events};
\node[svc]   (parser) at (9.5, 4.3) {parser\\\tiny:8008};
\node[svc]   (chunk)  at (9.5, 3.5) {chunking\\\tiny:8009};
\node[svc]   (embed)  at (9.5, 2.7) {embedding\\\tiny:8010};
\node[svc]   (idx)    at (9.5, 1.9) {indexing\\\tiny:8011};

% ── Support services (y = 1.0) ───────────────────────────────────
\node[supp] (admin_s) at (2.5, 0.5) {admin\\\tiny:8012};
\node[supp] (eval_s)  at (5.5, 0.5) {eval\\\tiny:8013};

% ── Infrastructure (y = -0.8) ────────────────────────────────────
\draw[rounded corners=6pt, fill=gray!5, draw=gray!40]
    (-0.2, -1.7) rectangle (12.7, -0.2);
\node[font=\footnotesize\bfseries, gray!60] at (6.25, -0.4) {Infrastructure};
\node[infra] (mongo)    at (2.0,  -1.1) {MongoDB\\\tiny:27017};
\node[infra] (weaviate) at (6.25, -1.1) {Weaviate\\\tiny:8080};
\node[infra] (rmq_i)    at (10.5, -1.1) {RabbitMQ\\\tiny:5672};

% ── External (y = -3.0) ──────────────────────────────────────────
\draw[rounded corners=6pt, fill=purple!3, draw=purple!30, dashed]
    (-0.2, -3.4) rectangle (12.7, -2.2);
\node[font=\footnotesize\bfseries, purple!50] at (6.25, -2.35) {External APIs};
\node[ext] (llm_e)  at (3.0,  -2.95) {Anthropic LLM\\{\tiny claude-sonnet-4-6}};
\node[ext] (emb_e)  at (9.0,  -2.95) {Embedding Provider\\{\tiny BGE-M3 / Cohere}};

% ── Arrows ───────────────────────────────────────────────────────
% clients → gateway
\draw[arr] (user)    -- (gw);
\draw[arr] (admin_u) -- (gw);

% gateway → auth (JWT validation)
\draw[arr, dashed] (gw.east) -- node[lbl,above=1pt]{JWT check} (auth.west);

% gateway → query pipeline
\draw[arr] (gw.south) -- node[lbl,right=1pt]{query} (orch.north);

% gateway → ingestion
\draw[arr] (gw.east) -- ++(0.6,0) -- ++(0,-0.9) -- node[lbl,above=1pt]{upload} (ingest.north);

% internal query pipeline (↓)
\draw[arr] (orch)   -- (qp);
\draw[arr] (qp)     -- (retr);
\draw[arr] (retr)   -- (rerank);
\draw[arr] (rerank) -- (gen_s);

% internal ingestion pipeline (↓)
\draw[arr] (ingest) -- (rmq_s);
\draw[arr] (rmq_s)  -- (parser);
\draw[arr] (parser) -- (chunk);
\draw[arr] (chunk)  -- (embed);
\draw[arr] (embed)  -- (idx);

% retrieval + indexing ↔ Weaviate
\draw[arr, <->] (retr.south) -- ++(0,-0.35) -- ++(3.75,0) -- (weaviate.north);
\draw[arr]      (idx.south)  -- ++(0,-0.35) -- ++(-3.25,0) -- (weaviate.north);

% generation → LLM external
\draw[arr] (gen_s.west)  -- ++(-0.8,0) -- ++(0,-3.2) -- (llm_e.north);

% embedding → embedding provider
\draw[arr] (embed.east) -- ++(0.8,0) -- ++(0,-2.5) -- (emb_e.north);

% eval ← RabbitMQ (query.completed)
\draw[arr, dashed] (rmq_i.west) -- ++(-4.75,0) -- (eval_s.south);

\end{tikzpicture}
\caption{Container diagram showing all 14 \medrag{} services grouped into the
  synchronous Query Pipeline (blue) and the asynchronous Ingestion Pipeline (orange).
  Dashed arrows indicate event-driven or optional communication.
  The eval service consumes \texttt{query.completed} events from RabbitMQ.}
\label{fig:c4_container}
\end{figure}

% ------------------------------------------------------------
\section{Query Pipeline}
\label{sec:query_pipeline}

% TODO: 3 pages
% Synchronous REST + SSE streaming
% Sequence diagram: gateway → auth → orchestrator → QP → retrieval → reranker → generation → SSE
% RAG mode switching in orchestrator: strategy pattern
% Conversation persistence in MongoDB

\begin{figure}[htbp]
\centering
\resizebox{\linewidth}{!}{%
\begin{tikzpicture}[
  part/.style = {rectangle, draw=blue!50, fill=blue!7, thick,
                 minimum width=2.8cm, minimum height=0.75cm,
                 align=center, font=\footnotesize\bfseries},
  msg/.style  = {-{Stealth[length=4pt]}, thick, black!65},
  ret/.style  = {-{Stealth[length=4pt]}, thick, black!40, dashed},
  lbl/.style  = {font=\scriptsize, fill=white, inner sep=1pt, align=center},
  act/.style  = {fill=blue!18, draw=blue!45, thin},
]

% participant boxes
\node[part] at (0,    0) {Client};
\node[part] at (3.5,  0) {API Gateway};
\node[part] at (7.5,  0) {Orchestrator};
\node[part] at (12.0, 0) {Retrieval\,+\,Reranker};
\node[part] at (16.5, 0) {Generation};

% lifelines
\foreach \x in {0, 3.5, 7.5, 12.0, 16.5}{
  \draw[gray!40, dashed, thin] (\x, -0.38) -- (\x, -9.8);
}

% 1. POST /chat/query
\draw[msg] (0,-1.0) -- node[lbl,above]{POST /chat/query} (3.5,-1.0);

% 2. forward (JWT validated)
\draw[msg] (3.5,-1.8) -- node[lbl,above]{validate JWT + forward} (7.5,-1.8);

% 3. [optional] activation box on orchestrator = query rewriting / HyDE
\draw[act] (7.38,-2.4) rectangle (7.62,-3.4);
\node[lbl, gray!70, align=left] at (9.0,-2.9) {\textit{[opt.]} rewrite query\\HyDE / query rewriting};

% 4. retrieve
\draw[msg] (7.5,-3.5) -- node[lbl,above]{retrieve($q$,\;$k{=}20$)} (12.0,-3.5);

% 5. activation on retrieval+reranker
\draw[act] (11.88,-3.5) rectangle (12.12,-4.6);

% 6. return top-k passages
\draw[ret] (12.0,-4.7) -- node[lbl,above]{top-5 passages} (7.5,-4.7);

% 7. generate
\draw[msg] (7.5,-5.5) -- node[lbl,above]{generate($q$,\;passages)} (16.5,-5.5);

% 8. activation on generation (LLM call)
\draw[act] (16.38,-5.5) rectangle (16.62,-7.2);
\node[lbl, gray!70] at (18.2,-6.35) {Anthropic\\LLM call};

% 9. SSE stream tokens → gateway → client
\draw[ret] (16.5,-7.3) -- node[lbl,above]{SSE stream} (3.5,-7.3);
\draw[ret] (3.5, -8.0) -- node[lbl,above]{SSE stream (proxied)} (0,-8.0);

% 10. async eval event
\draw[msg, dotted, thick] (7.5,-9.0) -- node[lbl,above]{\texttt{query.completed} (async)} (12.0,-9.0);

\end{tikzpicture}}
\caption{Query pipeline sequence. The API Gateway validates the JWT and forwards
  the request to the Orchestrator. Optional query rewriting or HyDE expansion
  (activation box) precedes retrieval. Generation streams tokens back to the client
  via Server-Sent Events. A \texttt{query.completed} event is published
  asynchronously for the eval service.}
\label{fig:seq_query}
\end{figure}

% ------------------------------------------------------------
\section{Ingestion Pipeline}
\label{sec:ingestion_pipeline}

% TODO: 2 pages
% Async, event-driven via RabbitMQ
% Steps: upload → parse → chunk → embed → index → status update
% Content-hash deduplication
% Chunking strategy comparison

\begin{figure}[htbp]
\centering
\begin{tikzpicture}[
  svc/.style   = {rectangle, rounded corners=5pt,
                  draw=blue!50, fill=blue!6, thick,
                  minimum width=2.6cm, minimum height=0.9cm, align=center,
                  font=\small},
  mq/.style    = {rectangle, rounded corners=5pt,
                  draw=orange!60, fill=orange!8, thick,
                  minimum width=2.6cm, minimum height=0.9cm, align=center,
                  font=\small},
  store/.style = {cylinder, draw=gray!60, fill=gray!9, thick,
                  shape border rotate=90, aspect=0.25,
                  minimum width=2.4cm, minimum height=0.85cm, align=center,
                  font=\small},
  arr/.style   = {-{Stealth[length=5pt]}, thick, black!65},
  sub/.style   = {font=\scriptsize, fill=white, inner sep=1pt},
]

% ── nodes (left to right) ─────────────────────────────────────────
\node[svc]   (ingest)  at ( 0,   0) {Ingestion\\:8007};
\node[mq]    (rmq)     at ( 3.4, 0) {RabbitMQ\\broker};
\node[svc]   (parser)  at ( 6.8, 0) {Parser\\:8008};
\node[svc]   (chunk)   at (10.2, 0) {Chunking\\:8009};
\node[svc]   (emb)     at (10.2,-2.0) {Embedding\\:8010};
\node[svc]   (idx)     at ( 6.8,-2.0) {Indexing\\:8011};

% ── storage nodes ─────────────────────────────────────────────────
\node[store] (mongo)    at ( 0,  -3.5) {MongoDB};
\node[store] (weaviate) at ( 6.8,-3.5) {Weaviate};

% ── event labels ──────────────────────────────────────────────────
\draw[arr] (ingest) -- node[sub,above=1pt]{\texttt{file.uploaded}} (rmq);
\draw[arr] (rmq)    -- node[sub,above=1pt]{\texttt{consume}} (parser);
\draw[arr] (parser) -- node[sub,above=1pt]{\texttt{file.parsed}} (chunk);
\draw[arr] (chunk.south) -- ++(0,-0.5) -- ++(3.4,0) -- (emb.north);
\node[sub] at (8.5,-0.65) {\texttt{chunks.created}};
\draw[arr] (emb.west) -- node[sub,above=1pt]{\texttt{embeddings.ready}} (idx.east);

% ── storage arrows ────────────────────────────────────────────────
\draw[arr] (ingest.south) -- node[sub,right=1pt]{dedup check} (mongo.north);
\draw[arr] (idx.south)    -- node[sub,right=1pt]{upsert vectors} (weaviate.north);
\draw[arr] (idx.west)     -- ++(-3.8,0) -- node[sub,right=1pt]{update status} (mongo.east);

\end{tikzpicture}
\caption{Ingestion pipeline: asynchronous event-driven document processing.
  Each stage publishes a completion event to RabbitMQ, consumed by the next
  stage. The Ingestion service performs content-hash deduplication against
  MongoDB before publishing. Final vectors are written to Weaviate by the
  Indexing service.}
\label{fig:seq_ingestion}
\end{figure}

% ------------------------------------------------------------
\section{Data Model}
\label{sec:data_model}

% TODO: 3 pages
% MongoDB collections: users, projects, documents, chunks, conversations, eval_results
% Weaviate schema: MedRAGChunk
% Rationale: MongoDB for metadata + Weaviate for vectors

\begin{figure}[htbp]
  \centering
  % TODO: TikZ UML class-style data model
  \fbox{\parbox{0.9\textwidth}{\centering\vspace{3cm}
    [Figure: MongoDB Data Model]\vspace{3cm}}}
  \caption{MongoDB collection schema. Each document links to its parent project;
    chunks carry a reference to their source document and strategy metadata.}
  \label{fig:data_model}
\end{figure}

% ------------------------------------------------------------
\section{Security Design}
\label{sec:security}

% TODO: 1 page
% JWT validation at gateway, role-based access (admin vs. user)
% Project-scoped retrieval: a user cannot query another project's documents
% bcrypt password hashing, configurable JWT expiry
